\chapter{Índices, Transações e Funções}
\label{cap:indices_transacoes_funcoes}

Esta aula encerra a Parte III deslocando o foco da SQL puramente estrutural e operacional para temas que aproximam o estudante do uso profissional de um SGBD: desempenho, consistência e encapsulamento de lógica. Até aqui, o leitor aprendeu a criar tabelas, manipular dados, elaborar consultas mais sofisticadas, reutilizar consultas por meio de views e controlar permissões com roles e privilégios. Agora, o objetivo é compreender três mecanismos centrais da prática administrativa e de desenvolvimento em banco de dados:

\begin{itemize}
    \item \textbf{índices}, que ajudam a melhorar o desempenho das consultas;
    \item \textbf{transações}, que preservam consistência em operações compostas;
    \item \textbf{funções}, que permitem encapsular lógica reutilizável dentro do banco.
\end{itemize}

Esses recursos mostram que um banco de dados não é apenas um repositório passivo de tabelas e linhas. Ele também é um ambiente de otimização, proteção e automação de processos.

\section{Índices}

\subsection{O que são índices}

Índices são estruturas auxiliares mantidas pelo SGBD para acelerar a localização de registros. Em termos intuitivos, um índice funciona de modo semelhante ao índice remissivo de um livro: em vez de percorrer todas as páginas até encontrar um assunto, consulta-se uma estrutura intermediária que aponta mais rapidamente para o local desejado.

Sem índice, muitas consultas dependem de varredura ampla da tabela. Com índice, o banco pode localizar registros de forma mais eficiente em situações apropriadas.

\subsection{Por que índices são importantes}

Índices tendem a ser úteis quando:

\begin{itemize}
    \item a tabela possui grande volume de dados;
    \item uma coluna é usada frequentemente em filtros;
    \item a ordenação por determinada coluna é recorrente;
    \item junções dependem repetidamente dos mesmos atributos;
    \item relatórios usam sempre os mesmos critérios de busca.
\end{itemize}

Em tabelas pequenas, o ganho pode ser pequeno. Em tabelas grandes, a diferença pode ser decisiva.

\subsection{Custo dos índices}

Embora acelerem muitas leituras, índices não são gratuitos. Eles:

\begin{itemize}
    \item ocupam espaço em disco;
    \item precisam ser mantidos em inserções, atualizações e remoções;
    \item podem aumentar o custo de operações de escrita;
    \item exigem planejamento para fazer sentido.
\end{itemize}

\begin{danger}
Criar índices em excesso ou sem critério pode degradar o desempenho geral do sistema, especialmente em ambientes com muita escrita.
\end{danger}

\section{Criação básica de índice}

\begin{sqlcode}
CREATE INDEX idx_produtos_nome ON produtos(nome);
\end{sqlcode}

Esse comando cria um índice sobre a coluna \texttt{nome} da tabela \texttt{produtos}. Consultas que filtram ou ordenam por essa coluna podem se beneficiar da estrutura.

\section{Tipos de índices no PostgreSQL}

O PostgreSQL oferece diferentes estratégias de indexação, cada uma mais adequada a certos tipos de consulta.

\subsection{B-tree}

O índice B-tree é o padrão do comando \texttt{CREATE INDEX}. É apropriado para:

\begin{itemize}
    \item igualdade;
    \item comparações de ordem;
    \item intervalos;
    \item ordenações.
\end{itemize}

\begin{sqlcode}
CREATE INDEX idx_aluno_nome ON alunos(nome);
\end{sqlcode}

\subsection{Hash}

O índice hash é voltado principalmente a comparações por igualdade.

\begin{sqlcode}
CREATE INDEX idx_clientes_email_hash
ON clientes USING hash (email);
\end{sqlcode}

\subsection{GIN}

O índice GIN é adequado para dados compostos, arrays e certos cenários de busca textual.

\begin{sqlcode}
CREATE EXTENSION IF NOT EXISTS pg_trgm;
CREATE INDEX idx_produtos_nome_trgm
ON produtos USING gin (nome gin_trgm_ops);
\end{sqlcode}

\begin{sqlcode}
SELECT produto_id, nome
FROM produtos
WHERE nome ILIKE '%caixa%';
\end{sqlcode}

\subsection{GIN com arrays}

\begin{sqlcode}
CREATE TABLE interesses (
    id SERIAL PRIMARY KEY,
    tags TEXT[]
);
CREATE INDEX idx_interesses_tags_gin
ON interesses USING gin (tags);
\end{sqlcode}

\begin{sqlcode}
SELECT *
FROM interesses
WHERE tags @> ARRAY['promoção', 'verão'];
\end{sqlcode}

\subsection{GiST}

O índice GiST é útil para cenários mais sofisticados, como intervalos ou dados espaciais.

\begin{sqlcode}
CREATE TABLE reservas (
    id SERIAL,
    periodo daterange
);
CREATE INDEX idx_reservas_gist
ON reservas USING gist (periodo);
\end{sqlcode}

\begin{sqlcode}
SELECT *
FROM reservas
WHERE periodo && daterange('2025-01-10', '2025-01-20');
\end{sqlcode}

\section{Índices multicoluna}

Índices podem ser criados com mais de uma coluna.

\begin{sqlcode}
CREATE INDEX idx_matricula_aluno_turma
ON matricula(aluno_id, turma_id);
\end{sqlcode}

Índices multicoluna são úteis quando consultas envolvem repetidamente a mesma combinação de atributos. Seu aproveitamento depende, entre outros fatores, da ordem das colunas no índice e da forma da consulta.

\section{Criação concorrente e reindexação}

\subsection{Criação concorrente}

\begin{sqlcode}
CREATE INDEX CONCURRENTLY idx_nome ON tabela(coluna);
\end{sqlcode}

A criação concorrente reduz bloqueios de escrita durante a construção do índice, sendo útil em produção.

\subsection{Reindexação}

\begin{sqlcode}
REINDEX TABLE nome_tabela;
\end{sqlcode}

A reindexação pode ser útil em cenários de fragmentação ou manutenção.

\section{Exemplo resolvido 1 -- Índice para busca frequente}

Suponha que a tabela \texttt{alunos} seja consultada com frequência pelo nome.

\begin{sqlcode}
CREATE INDEX idx_alunos_nome
ON alunos(nome);
\end{sqlcode}

\textbf{Interpretação:}
\begin{itemize}
    \item o índice tende a beneficiar buscas como \texttt{WHERE nome = ...};
    \item também pode ajudar em ordenações por nome;
    \item sua utilidade cresce com o tamanho da tabela e a frequência da consulta.
\end{itemize}

\section{Exemplo resolvido 2 -- Índice multicoluna para matrícula}

Suponha que o sistema acadêmico consulte constantemente matrículas por aluno e turma.

\begin{sqlcode}
CREATE INDEX idx_matricula_aluno_turma
ON matricula(aluno_id, turma_id);
\end{sqlcode}

\textbf{Interpretação:}
\begin{itemize}
    \item a combinação escolhida reflete um padrão de acesso esperado;
    \item o índice pode favorecer consultas conjuntas sobre esses dois campos;
    \item é mais útil do que criar índices sem relação com o comportamento real da aplicação.
\end{itemize}

\section{Transações}

\subsection{O que é uma transação}

Uma transação é uma sequência de operações tratada como unidade lógica de trabalho. Seu objetivo é garantir que um processo composto preserve a consistência do banco.

Em muitos casos, uma única operação não basta. É preciso executar duas ou mais ações de forma coordenada. Se apenas parte delas ocorrer, o banco pode ficar inconsistente.

\subsection{Exemplo intuitivo}

Pense em uma matrícula de aluno em turma com controle de vagas:
\begin{enumerate}
    \item inserir a matrícula;
    \item reduzir o número de vagas disponíveis.
\end{enumerate}

Se a matrícula for inserida e a vaga não for reduzida, o banco passa a representar uma realidade incoerente. Por isso, as duas ações devem compor a mesma transação.

\section{Propriedades ACID}

As transações são tradicionalmente descritas pelas propriedades ACID:

\subsection{Atomicidade}

Ou todas as operações da transação acontecem, ou nenhuma acontece.

\subsection{Consistência}

A transação deve levar o banco de um estado válido para outro estado válido.

\subsection{Isolamento}

Transações concorrentes não devem interferir indevidamente entre si.

\subsection{Durabilidade}

Uma vez confirmadas, as alterações persistem mesmo diante de falhas posteriores.

\section{Comandos principais de transação}

\subsection{BEGIN}

Inicia a transação.

\subsection{COMMIT}

Confirma definitivamente as alterações.

\subsection{ROLLBACK}

Desfaz o que foi feito desde o início da transação.

\section{Exemplo simples de transação}

\begin{sqlcode}
BEGIN;
UPDATE contas SET saldo = saldo - 100 WHERE conta_id = 1;
UPDATE contas SET saldo = saldo + 100 WHERE conta_id = 2;
COMMIT;
\end{sqlcode}

\textbf{Interpretação:}
\begin{itemize}
    \item a operação representa uma transferência;
    \item a consistência depende de ambas as atualizações;
    \item o \texttt{COMMIT} torna o processo definitivo.
\end{itemize}

\section{Exemplo com cancelamento}

\begin{sqlcode}
BEGIN;
UPDATE contas SET saldo = saldo - 100 WHERE conta_id = 1;
ROLLBACK;
\end{sqlcode}

Nesse caso, a operação é desfeita antes da confirmação final.

\section{Pontos de salvamento}

Em operações mais longas, podem ser usados \texttt{SAVEPOINT}s para desfazer apenas parte da transação.

\begin{sqlcode}
BEGIN;

UPDATE alunos
SET status = 'ATIVO'
WHERE aluno_id = 4;

SAVEPOINT ajuste_intermediario;

UPDATE alunos
SET email = 'diego.ramos@univ.edu'
WHERE aluno_id = 4;

ROLLBACK TO SAVEPOINT ajuste_intermediario;

COMMIT;
\end{sqlcode}

\section{Exemplo resolvido 3 -- Matrícula com redução de vaga}

Considere a rotina de matrícula em turma:

\begin{sqlcode}
BEGIN;

INSERT INTO matricula (aluno_id, turma_id)
VALUES (10, 3);

UPDATE turma
SET vagas_disponiveis = vagas_disponiveis - 1
WHERE turma_id = 3;

COMMIT;
\end{sqlcode}

\textbf{Interpretação:}
\begin{itemize}
    \item a matrícula e a redução de vagas formam um único processo;
    \item se uma das etapas falhar, o ideal é desfazer ambas;
    \item essa é exatamente a finalidade da transação.
\end{itemize}

\section{Exemplo resolvido 4 -- Cancelamento de matrícula}

\begin{sqlcode}
BEGIN;

DELETE FROM matricula
WHERE aluno_id = 10
  AND turma_id = 3;

UPDATE turma
SET vagas_disponiveis = vagas_disponiveis + 1
WHERE turma_id = 3;

COMMIT;
\end{sqlcode}

\textbf{Interpretação:}
\begin{itemize}
    \item cancelar matrícula e devolver vaga são partes do mesmo processo;
    \item a transação evita que apenas metade da operação se concretize.
\end{itemize}

\section{Funções}

\subsection{O que são funções}

Funções são blocos de código armazenados no servidor de banco de dados. Elas permitem encapsular lógica reutilizável que pode ser chamada a partir de consultas ou outras rotinas.

\subsection{Por que usar funções}

Funções ajudam a:
\begin{itemize}
    \item centralizar lógica repetitiva;
    \item reduzir duplicação de código SQL;
    \item organizar melhor o banco;
    \item tornar certas operações reutilizáveis;
    \item aproximar o banco de regras de negócio mais elaboradas.
\end{itemize}

\subsection{Características gerais}

Funções:
\begin{itemize}
    \item recebem parâmetros;
    \item retornam um valor ou conjunto de valores;
    \item executam no servidor;
    \item participam do contexto transacional;
    \item não devem executar \texttt{COMMIT} ou \texttt{ROLLBACK}.
\end{itemize}

\section{Estrutura geral de uma função}

\begin{sqlcode}
CREATE OR REPLACE FUNCTION nome_funcao(parametro tipo)
RETURNS tipo_retorno
LANGUAGE plpgsql
AS $$
BEGIN
    -- corpo da função
    RETURN valor;
END;
$$;
\end{sqlcode}

\section{Exemplo conceitual de função}

Criar uma função que recebe o identificador de um cliente e retorna o total gasto em pedidos:

\begin{sqlcode}
CREATE OR REPLACE FUNCTION total_gasto_cliente(p_cliente_id INTEGER)
RETURNS NUMERIC
LANGUAGE plpgsql
AS $$
DECLARE
    v_total NUMERIC;
BEGIN
    SELECT COALESCE(SUM(valor_total), 0)
    INTO v_total
    FROM pedidos
    WHERE cliente_id = p_cliente_id;

    RETURN v_total;
END;
$$;
\end{sqlcode}

\subsection{Chamada da função}

\begin{sqlcode}
SELECT total_gasto_cliente(3);
\end{sqlcode}

\section{Exemplo resolvido 5 -- Função para contar matrículas}

Suponha que se deseje uma rotina reutilizável para contar quantas matrículas um aluno possui.

\begin{sqlcode}
CREATE OR REPLACE FUNCTION total_matriculas_aluno(p_aluno_id INTEGER)
RETURNS INTEGER
LANGUAGE plpgsql
AS $$
DECLARE
    v_total INTEGER;
BEGIN
    SELECT COUNT(*)
    INTO v_total
    FROM matricula
    WHERE aluno_id = p_aluno_id;

    RETURN v_total;
END;
$$;
\end{sqlcode}

\begin{sqlcode}
SELECT total_matriculas_aluno(5);
\end{sqlcode}

\textbf{Interpretação:}
\begin{itemize}
    \item a lógica fica centralizada;
    \item outras consultas podem reutilizar a função;
    \item o banco passa a oferecer uma operação lógica pronta.
\end{itemize}

\section{Erro comum 1 -- Criar índice sem necessidade}

É um erro comum criar índices apenas por intuição, sem observar o padrão real de consultas.

\section{Erro comum 2 -- Supor que mais índices sempre significam melhor desempenho}

Índices ajudam leitura, mas aumentam o custo de escrita. Mais índices não significa automaticamente melhor banco.

\section{Erro comum 3 -- Não usar transação em operação composta}

Inserir matrícula e alterar vagas em comandos independentes, sem transação, é arriscado.

\section{Erro comum 4 -- Tentar usar COMMIT ou ROLLBACK dentro de função}

Funções participam do contexto transacional, mas não devem controlar diretamente a finalização da transação.

\section{Erro comum 5 -- Criar funções sem responsabilidade clara}

Funções muito genéricas ou excessivamente complexas dificultam manutenção.

\section{Boas práticas}

\begin{itemize}
    \item criar índices com base em consultas reais;
    \item revisar o impacto de índices em escrita;
    \item usar transações em processos compostos e críticos;
    \item nomear índices e funções de forma clara;
    \item manter funções focadas em uma responsabilidade bem definida;
    \item documentar quando um índice ou função atende a uma necessidade específica do sistema.
\end{itemize}

\begin{quadro_dica}
Em atividades práticas, vale manter um script de carga inicial do banco. Assim, testes com índices, transações e funções podem ser repetidos sem comprometer permanentemente o ambiente.
\end{quadro_dica}

\section{Minicaso integrador -- Secretaria acadêmica com desempenho e consistência}

Considere um sistema acadêmico com as tabelas \texttt{alunos}, \texttt{turma}, \texttt{matricula} e \texttt{notas}. A coordenação deseja:

\begin{enumerate}
    \item melhorar buscas por nome de aluno;
    \item acelerar consultas de matrícula por aluno e turma;
    \item garantir consistência no processo de matrícula;
    \item reutilizar a lógica de contagem de matrículas por aluno.
\end{enumerate}

Uma solução possível é:

\subsection*{Passo 1 -- Criar índice para busca por nome}

\begin{sqlcode}
CREATE INDEX idx_alunos_nome
ON alunos(nome);
\end{sqlcode}

\subsection*{Passo 2 -- Criar índice multicoluna}

\begin{sqlcode}
CREATE INDEX idx_matricula_aluno_turma
ON matricula(aluno_id, turma_id);
\end{sqlcode}

\subsection*{Passo 3 -- Proteger matrícula com transação}

\begin{sqlcode}
BEGIN;

INSERT INTO matricula (aluno_id, turma_id)
VALUES (12, 4);

UPDATE turma
SET vagas_disponiveis = vagas_disponiveis - 1
WHERE turma_id = 4;

COMMIT;
\end{sqlcode}

\subsection*{Passo 4 -- Criar função de contagem}

\begin{sqlcode}
CREATE OR REPLACE FUNCTION total_matriculas_aluno(p_aluno_id INTEGER)
RETURNS INTEGER
LANGUAGE plpgsql
AS $$
DECLARE
    v_total INTEGER;
BEGIN
    SELECT COUNT(*)
    INTO v_total
    FROM matricula
    WHERE aluno_id = p_aluno_id;

    RETURN v_total;
END;
$$;
\end{sqlcode}

\subsection*{Passo 5 -- Consultar a função}

\begin{sqlcode}
SELECT total_matriculas_aluno(12);
\end{sqlcode}

\textbf{Comentário didático:}
Esse minicaso mostra que a administração prática do banco envolve simultaneamente desempenho, consistência e reutilização lógica. A SQL passa a operar não apenas como linguagem de consulta, mas também como instrumento de engenharia do próprio banco.

\section{Exercícios básicos}

\begin{enumerate}
    \item Crie um índice B-tree para uma coluna consultada com frequência.
    \item Explique o objetivo de um índice.
    \item Escreva um exemplo com \texttt{BEGIN}, \texttt{COMMIT} e \texttt{ROLLBACK}.
    \item Defina, com suas palavras, as propriedades ACID.
    \item Crie uma função simples que receba um parâmetro e retorne um valor.
\end{enumerate}

\section{Exercícios intermediários}

\begin{enumerate}
    \item Compare índice B-tree e índice GIN.
    \item Explique em que situação um índice multicoluna pode ser útil.
    \item Crie uma transação que registre uma matrícula e reduza vagas disponíveis.
    \item Escreva uma função que conte quantos registros existem em uma tabela para um identificador informado.
    \item Explique por que funções não devem executar \texttt{COMMIT} ou \texttt{ROLLBACK}.
\end{enumerate}

\section{Exercícios avançados}

\begin{enumerate}
    \item Modele um cenário em que um índice GiST seja apropriado.
    \item Elabore uma transação com \texttt{SAVEPOINT} e desfazimento parcial.
    \item Projete uma função para calcular média de notas de um aluno.
    \item Explique como índices podem melhorar leitura e, ao mesmo tempo, aumentar custo de escrita.
\end{enumerate}

\section{Exercício integrador}

Considere um sistema acadêmico com as tabelas \texttt{alunos}, \texttt{turma}, \texttt{matricula}, \texttt{notas} e \texttt{professor}. Elabore:

\begin{enumerate}
    \item dois índices considerados úteis para o sistema;
    \item uma transação completa de matrícula;
    \item uma transação de cancelamento;
    \item uma função que retorne a quantidade de matrículas de um aluno;
    \item uma justificativa para o uso de cada recurso.
\end{enumerate}

\section{Síntese da aula}

Esta aula concluiu a Parte III mostrando que o uso profissional da SQL ultrapassa a criação de tabelas e a escrita de consultas. Foram abordados os índices como instrumentos de desempenho, as transações como mecanismos de consistência e as funções como recurso de encapsulamento lógico. Com exemplos resolvidos, erros comuns, exercícios progressivos e um minicaso integrador, o estudante encerra um percurso que vai da implementação estrutural à administração básica e à automação de rotinas no próprio banco de dados.